% ============================================================
% Appendix: formal statements and complete proofs (T1--P7)
% Spectral Submersion v3 -- ML4AL @ ACL submission
% Conventions: \Vx = sign vocabulary, G <= S_{|\Vx|} relabeling group.
% All citations marked [VERIFY] must be checked against the primary
% source before submission (Rule 4 of the editing protocol).
% ============================================================

\newcommand{\Vx}{\mathcal{V}_X}
\newcommand{\Orb}{\mathrm{Orb}}
\newcommand{\eff}{\mathrm{eff}}

\section{Formal framework and proofs}
\label{app:proofs}

\subsection{Setup and notation}
\label{app:setup}

Let $\Vx$ be a finite sign vocabulary, $|\Vx|=d$, and let a corpus
$X=(x_1,\dots,x_n)\in\Vx^{n}$ be generated by a probability kernel
$p(\cdot\mid\theta)$ indexed by a hypothesis $\theta\in\Theta$ (a
correspondence between signs and units of a candidate language, or any
richer latent object). A \emph{relabeling} is a permutation
$\sigma\in G\le S_{d}$ acting on corpora tokenwise,
$\sigma X=(\sigma x_1,\dots,\sigma x_n)$, and on hypotheses by
$(\sigma\cdot\theta)$, the hypothesis obtained by renaming each sign
$x$ to $\sigma(x)$ in the domain of $\theta$.

\begin{definition}[Relabeling equivariance]
\label{def:equivariance}
The model is \emph{$G$-equivariant} if for all $\sigma\in G$,
$\theta\in\Theta$ and corpora $X$,
\begin{equation}
p(\sigma X\mid \sigma\cdot\theta)=p(X\mid\theta).
\label{eq:equivariance}
\end{equation}
A statistic $T$ is \emph{$G$-invariant} if $T(\sigma X)=T(X)$ for all
$\sigma\in G$ and all $X$.
\end{definition}

Equation~\eqref{eq:equivariance} formalizes the epistemic situation of
an undeciphered script: nothing in the generative story changes if the
signs are consistently renamed. Frequencies of \emph{isomorphism
classes} of co-occurrence structures, spectra of PPMI matrices, degree
distributions, entropy curves, effective ranks, and any quantity
computed after quotienting by sign identity are $G$-invariant.
Raw co-occurrence \emph{matrices} are $G$-equivariant
($C(\sigma X)=P_\sigma C(X)P_\sigma^{\top}$); their spectra are
invariant.

\subsection{T1: non-identifiability as a group action}

\begin{theorem}[Orbit-constancy of the posterior]
\label{thm:T1}
Assume the model is $G$-equivariant and the prior $\pi$ is
$G$-invariant, $\pi(\sigma\cdot\theta)=\pi(\theta)$. Let $T$ be any
$G$-invariant statistic. Then for every $\sigma\in G$,
\[
p(\sigma\cdot\theta\mid T=t)\;=\;p(\theta\mid T=t)
\qquad\text{for all }t .
\]
Consequently the posterior is constant on each orbit
$\Orb_G(\theta)=\{\sigma\cdot\theta:\sigma\in G\}$.
\end{theorem}

\begin{proof}
Fix $t$ and $\sigma\in G$. The likelihood of $t$ under
$\sigma\cdot\theta$ is
\[
p(T=t\mid\sigma\cdot\theta)
=\sum_{X:\,T(X)=t} p(X\mid\sigma\cdot\theta).
\]
Substituting $X=\sigma X'$ (a bijection of $\Vx^{n}$) and using
\eqref{eq:equivariance},
\[
p(T=t\mid\sigma\cdot\theta)
=\sum_{X':\,T(\sigma X')=t} p(X'\mid\theta)
=\sum_{X':\,T(X')=t} p(X'\mid\theta),
\]
where the last equality uses $G$-invariance of $T$. Hence
$p(T=t\mid\sigma\cdot\theta)=p(T=t\mid\theta)$. By Bayes' rule and
$G$-invariance of $\pi$,
\[
p(\sigma\cdot\theta\mid t)
=\frac{\pi(\sigma\cdot\theta)\,p(t\mid\sigma\cdot\theta)}
       {p(t)}
=\frac{\pi(\theta)\,p(t\mid\theta)}{p(t)}
=p(\theta\mid t).\qedhere
\]
\end{proof}

\begin{corollary}[Data processing / no post-hoc identifiability]
\label{cor:T1-dpi}
Let $f$ be any measurable function and $T$ any $G$-invariant
statistic. Then $f\circ T$ is $G$-invariant, and
Theorem~\ref{thm:T1} applies verbatim to $f\circ T$. In particular, no
pipeline stage applied downstream of $G$-invariant statistics
(embeddings-up-to-rotation, spectra, clusterings of isomorphism
classes, neural post-processing) can break orbit-constancy of the
posterior.
\end{corollary}

\begin{proof}
$(f\circ T)(\sigma X)=f(T(\sigma X))=f(T(X))=(f\circ T)(X)$.
\end{proof}

\begin{remark}
Identifiability therefore requires external anchoring knowledge $K$
whose likelihood is \emph{not} $G$-invariant (iconographic referents,
documented reading traditions, bilinguals). This is the formal license
for the claim-level protocol of Definition~\ref{def:claims}.
\end{remark}

\subsection{T2: quantified impossibility (Fano)}

\begin{theorem}[Fano bound on any invariant decoder]
\label{thm:T2}
Let $\theta$ be uniform on an orbit (or any indistinguishability
class) $\mathcal{N}$ with $|\mathcal{N}|=N\ge 2$, let $T$ be a
$G$-invariant statistic of the corpus, and let
$\hat\theta=f(T)$ be any (possibly randomized) estimator. Then
\[
P_e \;=\;\Pr[\hat\theta\neq\theta]\;\ge\;
1-\frac{I(\theta;T)+\log 2}{\log N}
\;=\;1-\frac{\log 2}{\log N},
\]
since $I(\theta;T)=0$.
\end{theorem}

\begin{proof}
By Theorem~\ref{thm:T1} the likelihood $p(T=t\mid\theta)$ is the same
for all $\theta$ in the orbit; with $\theta$ uniform on the orbit this
means $T\perp\theta$, i.e.\ $I(\theta;T)=0$. Fano's inequality
[VERIFY: Cover \& Thomas, Thm.~2.10.1] gives
$H(\theta\mid T)\le h(P_e)+P_e\log(N-1)$, and
$H(\theta\mid T)=H(\theta)=\log N$. Bounding $h(P_e)\le\log 2$ and
$\log(N-1)\le\log N$ and rearranging yields the claim.
\end{proof}

\begin{definition}[$\varepsilon(n)$-empirical orbit]
\label{def:empirical-orbit}
Fix a corpus $X$ with $n$ tokens and the co-occurrence statistic
$\hat C$ with concentration radius $\varepsilon(n)$ from
Theorem~\ref{thm:T3}. The \emph{$\varepsilon(n)$-empirical orbit} of
$\theta$ is
\[
\Orb_{\varepsilon(n)}(\theta)=
\bigl\{\sigma\cdot\theta:\;
\|P_\sigma \mathbb{E}\hat C P_\sigma^{\top}-\mathbb{E}\hat C\|_2
\le 2\varepsilon(n)\bigr\},
\]
the set of relabelings whose \emph{population} statistics differ by
less than twice the sampling error, hence are statistically
indistinguishable at sample size $n$.
\end{definition}

\begin{corollary}[Finite-sample impossibility and the price of $K$]
\label{cor:T2-price}
With $N(n)=|\Orb_{\varepsilon(n)}(\theta)|$, any estimator based on
$\hat C$ satisfies
$P_e\ge 1-\bigl(o(1)+\log 2\bigr)/\log N(n)$.
Moreover, for a decoder additionally consuming external knowledge $K$
to achieve $P_e\le\delta$, Fano applied to the pair $(T,K)$ requires
\[
I(\theta;K)\;\ge\;(1-\delta)\log N(n)-\log 2 .
\]
That is: the minimum number of nats of genuinely anchoring information
scales with the log-size of the empirical indistinguishability class.
\end{corollary}

\begin{proof}
The first claim follows from Theorem~\ref{thm:T2} applied to the class
$\Orb_{\varepsilon(n)}(\theta)$, within which the distributions of
$\hat C$ overlap up to total-variation $o(1)$ by
Theorem~\ref{thm:T3} and Pinsker's inequality [VERIFY]. The second
claim is Fano with side information:
$H(\theta\mid T,K)\ge H(\theta)-I(\theta;T)-I(\theta;K)
=\log N(n)-I(\theta;K)$, combined with
$H(\theta\mid T,K)\le h(\delta)+\delta\log N(n)$.
\end{proof}

\begin{remark}
$N(n)$ is computable in practice by sampling permutations that
preserve the empirical frequency profile within tolerance
$\varepsilon(n)$; \texttt{[CALCULAR: script
\texttt{estimate\_empirical\_orbit.py}, entrada: corpus RR A--F,
n=5460]}.
\end{remark}

\subsection{T3: concentration of co-occurrence and PPMI}

\emph{Assumption (L).} The corpus consists of $m$ lines drawn i.i.d.\
from a line distribution, each line of length at most $L$ tokens;
co-occurrence windows do not cross line boundaries (as enforced by the
pipeline). This matches the tablet data model. An extension to
$w$-dependent token processes via block partitioning is sketched at
the end.

\begin{theorem}[Matrix Bernstein for windowed co-occurrence]
\label{thm:T3}
Let $\hat C=\frac1m\sum_{j=1}^{m}C_j$ where $C_j\in\mathbb{R}^{d\times d}$
is the (symmetrized) window-$w$ co-occurrence matrix of line $j$.
Then $\|C_j\|_2\le 2wL=:B$ and, with
$v=\|\mathbb{E}[(C_j-\mathbb{E}C_j)^2]\|_2\le B^2$, for all $t\ge0$
\[
\Pr\!\left[\|\hat C-\mathbb{E}\hat C\|_2\ge t\right]
\le 2d\exp\!\left(\frac{-mt^{2}/2}{v+Bt/3}\right).
\]
In particular, with probability at least $1-2d^{-1}$,
\[
\|\hat C-\mathbb{E}\hat C\|_2
\;\le\;
\sqrt{\frac{4v\log d}{m}}+\frac{4B\log d}{3m}
\;=\;O\!\left(wL\sqrt{\frac{\log d}{m}}\right)
=:\varepsilon(n),
\]
where $n\le mL$ is the token count.
\end{theorem}

\begin{proof}
Each line contributes at most $wL$ ordered co-occurring pairs; each
pair adds a matrix of spectral norm at most $2$ (a symmetrized
elementary matrix), whence $\|C_j\|_2\le 2wL$. The summands
$C_j-\mathbb{E}C_j$ are i.i.d., mean zero, bounded by $2B$; the matrix
Bernstein inequality [VERIFY: Tropp 2015, Thm.~1.6.2] applied to the
sum $\frac1m\sum_j (C_j-\mathbb{E}C_j)$ gives the tail bound; setting
the right side to $2d^{-1}$ and solving for $t$ gives the stated
radius.
\end{proof}

\begin{lemma}[PPMI perturbation under clipping]
\label{lem:ppmi}
Let $M_{ij}=\max\{0,\log\frac{p_{ij}}{p_i p_j}\}$ be computed with the
pipeline's clipping floor $p_{ij}\vee\epsilon$, $p_i\vee\epsilon$
(flag \texttt{pmi\_epsilon}). Then entrywise
\[
|\hat M_{ij}-M_{ij}|
\le \frac{|\hat p_{ij}-p_{ij}|}{\epsilon}
+\frac{|\hat p_i-p_i|}{\epsilon}
+\frac{|\hat p_j-p_j|}{\epsilon},
\]
hence
$\|\hat M-M\|_F\le \frac{3d}{\epsilon}\,
\max\{\|\hat p-p\|_\infty\text{-type terms}\}
=O\!\bigl(\tfrac{d}{\epsilon}\varepsilon(n)\bigr)$.
\end{lemma}

\begin{proof}
$x\mapsto\log(x\vee\epsilon)$ is $(1/\epsilon)$-Lipschitz on
$[0,1]$; $\max\{0,\cdot\}$ is $1$-Lipschitz; the triangle inequality
over the three estimated quantities gives the entrywise bound, and
summing $d^2$ entries gives the Frobenius bound.
\end{proof}

\begin{corollary}[Stability of spectral embeddings]
\label{cor:daviskahan}
Let $\delta_k=\sigma_k(M)-\sigma_{k+1}(M)$ be the singular gap. If
$\|\hat M-M\|_2\le\delta_k/4$, then the top-$k$ left singular
subspaces satisfy (Wedin's $\sin\Theta$ theorem [VERIFY: Wedin 1972;
Davis--Kahan 1970])
\[
\|\sin\Theta(\hat U_k,U_k)\|_2
\;\le\;\frac{2\,\|\hat M-M\|_2}{\delta_k}.
\]
Consequently the embedding geometry is stable whenever the spectral
gap dominates the concentration radius of
Theorem~\ref{thm:T3}--Lemma~\ref{lem:ppmi}.
\end{corollary}

\begin{remark}[Bootstrap as corollary, not observation]
The empirically observed bootstrap stability of the effective rank
(CV $=0.65\%$) and of embedding cosines ($0.77$) on the structured
synthetic corpus is the predicted behaviour of
Corollary~\ref{cor:daviskahan} in the regime
$\varepsilon(n)\ll\delta_k$, combined with
Lemma~\ref{lem:L5} below for $r_{\eff}$. The failed sanity checks on
short-line/large-vocabulary corpora (Indus real) correspond to the
opposite regime $\varepsilon(n)\gtrsim\delta_k$, where no such
stability is guaranteed --- the theory predicts both the successes and
the reported failures.
\end{remark}

\emph{Extension to $w$-dependence.} If lines cannot be assumed
i.i.d., partition the token sequence into $w$ interleaved blocks;
within each block the windowed contributions are independent, and a
union bound over the $w$ blocks multiplies the tail by $w$ and the
radius by $\sqrt{w}$, preserving
$\varepsilon(n)=O(\sqrt{w^{2}\log d/n})$ up to constants
[VERIFY: standard $m$-dependent Bernstein argument].

\subsection{T4: Procrustes recovery with noisy anchors}

\begin{theorem}[Perturbation of the Procrustes rotation]
\label{thm:T4}
Let $B\in\mathbb{R}^{k\times r}$ be the anchor matrix
($k$ anchors, dimension $r$, $k\ge r$), and suppose
$A=B\Omega^{*}+E$ with $\Omega^{*}\in O(r)$. Let
$\hat\Omega=UV^{\top}$ where $B^{\top}A=U\Sigma V^{\top}$. Then
\[
\|\hat\Omega-\Omega^{*}\|_F
\;\le\;\frac{2\,\|B^{\top}E\|_F}{\sigma_{r}(B)^{2}}
\;\le\;\frac{2\,\|B\|_2\,\|E\|_F}{\sigma_{r}(B)^{2}} .
\]
\end{theorem}

\begin{proof}
Write $M=B^{\top}B\,\Omega^{*}$ and $\hat M=B^{\top}A=M+B^{\top}E$.
Since $B^{\top}B$ is symmetric positive semi-definite, the polar
decomposition of $M$ has orthogonal factor exactly $\Omega^{*}$
(indeed $M=(B^{\top}B)\Omega^{*}$ with $B^{\top}B\succeq0$ is already
a polar decomposition read right-to-left), and
$\sigma_{r}(M)=\sigma_{r}(B^{\top}B)=\sigma_{r}(B)^{2}$. The
orthogonal Procrustes solution $\hat\Omega$ is the orthogonal polar
factor of $\hat M$. The perturbation bound for polar factors
[VERIFY: R.-C. Li 1995, ``New perturbation bounds for the unitary
polar factor'']
\[
\|\hat\Omega-\Omega^{*}\|_F
\le\frac{2}{\sigma_r(M)+\sigma_r(\hat M)}\,\|\hat M-M\|_F
\le\frac{2\|B^{\top}E\|_F}{\sigma_r(B)^2}
\]
gives the claim (dropping the nonnegative $\sigma_r(\hat M)$ in the
denominator).
\end{proof}

\begin{proposition}[Margin condition for nearest-neighbour matching]
\label{prop:margin}
Let a test pair $(x,y)$ satisfy $y=x\Omega^{*}+e$, and define the
\emph{margin}
$\gamma(y)=\min_{y'\neq y^{\mathrm{match}}}\|y-y'\|
-\|y-y^{\mathrm{match}}\|$ over candidate targets. If
\[
\gamma(y)\;>\;2\bigl(\|x\|_2\,\|\hat\Omega-\Omega^{*}\|_2
+\|e\|_2\bigr),
\]
then nearest-neighbour matching under $\hat\Omega$ returns the correct
target for $x$.
\end{proposition}

\begin{proof}
$\|x\hat\Omega-y^{\mathrm{match}}\|
\le\|x\hat\Omega-x\Omega^{*}\|+\|x\Omega^{*}-y^{\mathrm{match}}\|
\le\|x\|\|\hat\Omega-\Omega^{*}\|+\|e\|$. For any wrong candidate
$y'$,
$\|x\hat\Omega-y'\|\ge\|y-y'\|-\|y-x\hat\Omega\|
\ge\|y-y^{\mathrm{match}}\|+\gamma
-\bigl(\|x\|\|\hat\Omega-\Omega^{*}\|+\|e\|\bigr)-\|e\|$~%
\footnote{Using $\|y-x\hat\Omega\|\le\|e\|+\|x\|\|\hat\Omega-\Omega^*\|$.}.
The stated margin condition makes the wrong-candidate distance
strictly exceed the correct-candidate distance.
\end{proof}

\begin{remark}[Theorem--experiment cross-validation: negative result]
We ran the planned cross-validation
(\texttt{fig\_margin\_distribution.py}, seed 42): on the synthetic
anchor benchmark the sufficient threshold
$2(\|x\|\|\hat\Omega-\Omega^{*}\|+\|e\|)$, instantiated with the
computable bound of Theorem~\ref{thm:T4} for
$\|\hat\Omega-\Omega^{*}\|$, is \emph{vacuous}: it certifies $0\%$ of
test pairs while the observed Acc@1 is $0.689$. Moreover the margin
statistic $\gamma$ as defined is only weakly predictive of
nearest-neighbour correctness on this benchmark (AUC $=0.536$,
$n=90$). We report this as a negative result: the bound is sufficient
but far from necessary here (the residual-based rotation-error bound
is loose when anchors are noisy and the true map is not an exact
isometry), and the benchmark's errors are governed by local crowding
rather than by the global margin. Tightening the constant (e.g.\ via
a direct estimate of $\|\hat\Omega-\Omega^{*}\|$ under a known
ground-truth rotation) is stated as future work, not claimed.
\end{remark}

\subsection{L5: Lipschitz stability of the effective rank}

\begin{lemma}[Effective rank is stable]
\label{lem:L5}
Let $\sigma,\sigma'\in\mathbb{R}^{k}_{\ge0}\setminus\{0\}$ with
normalizations $p=\sigma/\|\sigma\|_1$, $p'=\sigma'/\|\sigma'\|_1$,
and let $r_{\eff}(\sigma)=e^{H(p)}$. Then, with
$T=\tfrac12\|p-p'\|_1\le\tfrac{2\|\sigma-\sigma'\|_1}{\|\sigma\|_1}$:
\[
|r_{\eff}(\sigma)-r_{\eff}(\sigma')|
\;\le\;k\,\bigl(T\log(k-1)+h(T)\bigr),
\]
where $h$ is the binary entropy, by the Fannes--Audenaert inequality
[VERIFY: Audenaert 2007]. In particular $r_{\eff}$ is continuous, and
locally Lipschitz away from $T\to\tfrac12$ with modulus
$O(k\log k)\cdot\|\sigma-\sigma'\|_1/\|\sigma\|_1$.
\end{lemma}

\begin{proof}
$|e^{H(p)}-e^{H(p')}|\le e^{\max H}\,|H(p)-H(p')|\le k\,|H(p)-H(p')|$
by the mean value theorem and $H\le\log k$. Audenaert's sharpening of
Fannes' inequality bounds $|H(p)-H(p')|\le T\log(k-1)+h(T)$. The
normalization map satisfies
$\|p-p'\|_1\le\frac{2}{\|\sigma\|_1}\|\sigma-\sigma'\|_1$ by the
standard triangle-inequality computation
$\|\frac{\sigma}{\|\sigma\|_1}-\frac{\sigma'}{\|\sigma'\|_1}\|_1
\le\frac{\|\sigma-\sigma'\|_1}{\|\sigma\|_1}
+\bigl|\frac{1}{\|\sigma\|_1}\|\sigma'\|_1-1\bigr|
\le\frac{2\|\sigma-\sigma'\|_1}{\|\sigma\|_1}$.
\end{proof}

\begin{remark}
Combined with Theorem~\ref{thm:T3} and Weyl's inequality
($|\sigma_i(\hat M)-\sigma_i(M)|\le\|\hat M-M\|_2$), this yields a
complete chain: sampling noise $\to$ spectrum $\to$ effective rank,
justifying $r_{\eff}$ as the pipeline's primary robust structure
statistic (claim level C1).
\end{remark}

\subsection{P6: selection bias of fusion decoding}

\begin{proposition}[Monotonicity and the $7.31<8.26$ artefact]
\label{prop:P6}
Fix an input $x$ and a \emph{fixed} finite candidate set
$\mathcal{C}$ of output sequences (e.g.\ the union of beam candidate
pools). For $\lambda\in[0,1]$ define
\[
\hat y_\lambda=\arg\max_{y\in\mathcal{C}}\;
(1-\lambda)\log p_{\mathrm{model}}(y\mid x)
+\lambda\log p_{\mathrm{LM}}(y).
\]
Then:
\begin{enumerate}[label=(\roman*),nosep]
\item $\lambda\mapsto\log p_{\mathrm{LM}}(\hat y_\lambda)$ is
non-decreasing; equivalently the LM cross-entropy of the decoded
output is non-increasing in $\lambda$.
\item If $\mathcal{C}$ contains at least one sequence $y_0$ drawn
from any reference process $P_0$, then at $\lambda=1$,
$-\log p_{\mathrm{LM}}(\hat y_1)\le-\log p_{\mathrm{LM}}(y_0)$, so in
expectation over draws,
$\mathbb{E}\bigl[-\log p_{\mathrm{LM}}(\hat y_1)\bigr]\le
H_{\times}(P_0,p_{\mathrm{LM}})$, the cross-entropy of the real
process.
\item Consequently, comparing the LM bits-per-glyph of a
score-maximizing decoder to the \emph{mean} bits-per-glyph of real
held-out lines is structurally biased downward: an output below the
real-text cross-entropy (here $7.31<8.26$) is the expected behaviour
of the design, and must not be read as evidence that the generator is
``more authentic than the tablets''.
\end{enumerate}
\end{proposition}

\begin{proof}
(i) Let $\lambda_1<\lambda_2$ with optima $\hat y_1,\hat y_2$ and
write $m(y)=\log p_{\mathrm{model}}(y\mid x)$,
$\ell(y)=\log p_{\mathrm{LM}}(y)$. Optimality gives
\[
(1-\lambda_1)m(\hat y_1)+\lambda_1\ell(\hat y_1)\ge
(1-\lambda_1)m(\hat y_2)+\lambda_1\ell(\hat y_2),
\]
\[
(1-\lambda_2)m(\hat y_2)+\lambda_2\ell(\hat y_2)\ge
(1-\lambda_2)m(\hat y_1)+\lambda_2\ell(\hat y_1).
\]
Adding the two inequalities and cancelling yields
$(\lambda_2-\lambda_1)\bigl(\ell(\hat y_2)-\ell(\hat y_1)\bigr)\ge0$,
hence $\ell(\hat y_2)\ge\ell(\hat y_1)$.
(ii) Immediate from $\hat y_1=\arg\max_{\mathcal C}\ell$ and
$y_0\in\mathcal{C}$; take expectations.
(iii) By (i), for any $\lambda\in(0,1]$ the decoded output's LM
log-likelihood is at least that of the $\lambda=0$ output and is
pushed toward the maximum over $\mathcal C$; the comparison target
$8.26$ is a \emph{mean} over typical realizations, whereas the
decoder reports a \emph{maximum-selected} statistic. A mode or
near-mode of a distribution always has per-symbol surprisal at most
the mean surprisal plus $o(1)$ under coverage; the inequality
$7.31<8.26$ is therefore a selection effect, not an empirical
discovery about authenticity.
\end{proof}

\begin{remark}[Honest scope]
In the implementation the beam's candidate pool itself depends on
$\lambda$ through pruning, so (i) holds exactly for exact decoding or
a shared candidate pool, and approximately for beam search; the
$\lambda$-sweep experiment (Task~2.3) verifies the monotone trend
empirically. \texttt{[CALCULAR: script
\texttt{sweep\_lambda.py}, $\lambda\in\{0,0.05,\dots,1\}$, seed 42]}.
\end{remark}

\subsection{D7 + P7: the claim-level protocol, formalized}

\begin{definition}[Claim classes]
\label{def:claims}
Let $\mathcal{T}$ be the algebra of $G$-invariant statistics of the
corpus, and let $\mathcal{E}_0$ be a specified family of internal
null ensembles (permutation, frequency-matched, uniform resampling).
A scientific statement $\varphi$ about the script is classified:
\begin{itemize}[nosep]
\item \textbf{C1 (structural signal).} $\varphi$ is a statement about
the value of some $T\in\mathcal{T}$ whose distribution under the data
differs from its distribution under every ensemble in
$\mathcal{E}_0$; C1 statements are falsifiable by internal negative
controls (permutation tests).
\item \textbf{C2 (rankable hypothesis).} $\varphi$ asserts a
preference ordering over hypotheses induced by a $G$-invariant score
$s:\Theta\to\mathbb{R}$ with $s(\sigma\cdot\theta)=s(\theta)$; C2
statements select orbits (or distributions over orbits) and are
rankeable but not falsifiable from internal data alone.
\item \textbf{C3 (anchored semantics).} $\varphi$'s truth value is
not constant on some orbit $\Orb_G(\theta)$ --- it distinguishes
points \emph{within} an orbit (e.g.\ ``glyph 600 means
`frigatebird'\,'').
\end{itemize}
\end{definition}

\begin{proposition}[C3 is unreachable without anchors]
\label{prop:P7}
No decision procedure $\varphi=f(T)$ with $T\in\mathcal{T}$ decides a
C3 statement: formally,
\[
\text{C3}\;\cap\;
\{\text{statements decidable from } \mathcal{T}\text{ alone}\}
=\emptyset .
\]
\end{proposition}

\begin{proof}
Suppose $\varphi=f(T)$ decides a C3 statement. By
Corollary~\ref{cor:T1-dpi}, $f\circ T$ is $G$-invariant, so
$\varphi$'s value is constant on every orbit. But by
Definition~\ref{def:claims}, a C3 statement is non-constant on some
orbit --- contradiction.
\end{proof}

\begin{remark}[The protocol as a theorem, not a style guide]
Proposition~\ref{prop:P7} upgrades the C1/C2/C3 protocol from
editorial best practice to a formal consequence of
Theorem~\ref{thm:T1}: the three levels partition scientific
statements by \emph{what can falsify them} (internal nulls; nothing
internal; external anchors $K$), and the boundary between C2 and C3
is exactly the boundary drawn by the group action. This is the
paper's central conceptual contribution and is stated as such in the
introduction.
\end{remark}
